\documentclass[a4paper,11pt]{article}
\usepackage[english]{babel}
\usepackage{natbib}
\usepackage{url}
\usepackage[utf8x]{inputenc}
\usepackage{amsmath}
\usepackage{graphicx}
%\graphicspath{{img/}}
\usepackage{parskip}
\usepackage{subfig}
\usepackage{float}
\usepackage{tabularx}
\usepackage{fancyhdr}
\usepackage{parskip}
\usepackage{indentfirst}
\usepackage{color}
\usepackage[dvipsnames]{xcolor}
\setlength{\parindent}{15pt}
\usepackage{vmargin}
\usepackage{framed}
\usepackage[hidelinks]{hyperref}
%\setmarginsrb{3 cm}{2.5 cm}{3 cm}{2.5 cm}{1 cm}{1.5 cm}{1 cm}{1.5 cm}%

\usepackage{hyperref}

\usepackage{xcolor}
\hypersetup{
    colorlinks,
    linkcolor={red!50!black},
    citecolor={blue!50!black},
    urlcolor={blue!80!black}
}

\title{
Group 0 - Final Report
\rule{\linewidth}{0.2 mm} \\[0.4 cm]
\Large{\textbf{INFO085 - Compiler}} \\
Develop a compiler for the VSOP language
\rule{\linewidth}{0.2 mm} \\[0.4 cm]
\vspace*{8 cm}
}


\date{\today}

\author{Thibaut \textsc{Chavet} \and Maxim \textsc{Henry} \and Maxime \textsc{Massart}}



\begin{document}

\begin{figure}
\begin{center}
\includegraphics[scale = 1]{logoulg}
\end{center}
\end{figure}

\maketitle
\setcounter{page}{0}
\thispagestyle{empty}


\tableofcontents
\thispagestyle{empty}
\setcounter{page}{0}
\newpage

% %%%%%%%%%%%%%%%%%%%%%%%%%%%%%%%%%%%%%%%%%%%%%%%%%%%%%%%%%%%%%%%%%%%%%%%%%%%%%%%%%%%%%%%%%
% 
% \begin{abstract} 
% 
% \end{abstract}
% 
% 
% %%%%%%%%%%%%%%%%%%%%%%%%%%%%%%%%%%%%%%%%%%%%%%%%%%%%%%%%%%%%%%%%%%%%%%%%%%%%%%%%%%%%%%%%%
% 
% \tableofcontents
% \pagebreak
% 
% %%%%%%%%%%%%%%%%%%%%%%%%%%%%%%%%%%%%%%%%%%%%%%%%%%%%%%%%%%%%%%%%%%%%%%%%%%%%%%%%%%%%%%%%%

\section{Introduction}
This project is done for the compiler course given by Prof. Geurts during the academic year 2016-2017. 
We are asked to implement a vsop compiler. The vsop language is defined on the course website.

As defined in the assignment, \texttt{The Very Simple Object-oriented Programming} language (or VSOP) is, 
as its name implies, a simple object-oriented language.



\section{Our compiler}
We chose the C++ 11 language to implement our vsop compiler  because it is well known by the team and allows us
to use powerful tools such as \href{http://www.gnu.org/software/bison/}{GNU Bison} and 
\href{https://github.com/westes/flex}{GNU Flex}.

The implementation was divided into 4 main steps: 
\begin{itemize}
 \item the Lexical Analysis (cf \ref{sec:lexana})
 \item the Syntax Analysis (cf \ref{sec:syntax})
 \item the Semantic Checking (cf \ref{sec:semantic})
 \item the Code Generation (cf \ref{sec:codeGeneration}) 
\end{itemize}



\subsection*{Status of the compiler}

Our compiler computes the LLVM which is then being compiled to machine code using llvm compiler version 3.5.

\subsection{Lexical Analysis}
  \label{sec:lexana}
  The lexical analysis is used to make sure the words inside the code to be compiled are allowed by the language, parse the vsop file.
  The goal of the scanner we are then able to generate, is to create tokens to be translated into tree nodes in the next steps. 
  We generate the tokens using the \href{https://github.com/westes/flex}{GNU Flex} program. Flex will generate a .c and a .h files 
  which will be used to compile our complete compiler.

  From the moment, we implemented the syntax analysis, the yylex() function is called from the yyparse() which is used in the second step.
  It has the drawback of not allowing the simple tests performed in the first step as both errors coming from Lexical and Syntax analysis
  will be displayed as the file is analyzed. 

  We had an issue with this lexical part. As our C++ Bison yyparse() calls yylex itself, when the input of this part is not correct 
  in the vsop formalism, all the tokens are not read until the end, but only until the syntax error. 
  %We could have a first parsing phase that would be done twice but for performances, it is not a good option. 
  We thought about a simple solution which was to do the first phase(GNU FLex) until the end. Then launch do the second phase(GNU Bison) 
  which would also redo the first one. It is not that efficient after all but it would be useful to pass your automated tests.
  



  \subsubsection{Tokens}
    The tokens that we used are explained in Table \ref{tab:listToken}.

    Those tokens will be used by \href{http://www.gnu.org/software/bison/}{GNU Bison}.

    \begin{table}[hp]
      \centering
      \begin{tabular}{|r|p{10cm}|}
      \hline
      Token & Description\\\hline
      INT\_LIT & Represents a integer literal\\
      AND & Represents and boolean operation\\
      BOOL &Represents the type definition of boolean \\
      CLASS & Represents the class definition keyword\\
      DO & Represents the beginning of a while repeated routine\\
      ELSE & Represents the else keyword in the if statement\\
      EXTENDS & Represents the inheritance performed in class definition\\
      FALSE & Represents the false boolean literal \\
      IF & Represents the if statement keyword\\
      IN & Represents the in keyword used in  the let statement\\
      INT32 & Represents the type definition of an integer on 32 bits\\
      ISNULL & Represents the isnull conditional keyword\\
      LET & Represents the definition of local variables\\
      NEW & Represents the object instantiation keyword\\
      NOT & Represents the negation in conditions keyword\\
      STRING & Represents the keyword defining a string element\\
      THEN & Represents the keyword introducing the affirmative part of the condition contained in the if statement\\
      TRUE & Represents the true boolean literal\\
      UNIT & Represents the unit type (similar to void in C)\\
      WHILE & Represents the loop keyword\\
      $<$strval$>$ TYPE\_ID & Represents a Type identifier \\
      $<$strval$>$ OBJECT\_ID & Represents an object identifier\\
      $<$strval$>$ STRING\_LIT & Represents a string literal \\
      LBRACE & Represents the beginning of a block\\
      RBRACE& Represents the end of a block\\
      LPAR & Represents the beginning of a parenthesis used in function calls and functions declaration\\
      RPAR & Represents the end of a parenthesis\\
      COLON & Represents the colon symbols delimiting type definitions\\
      SEMICOLON & Represents the semi colon symbol at the end of a non-returned function instruction \\
      COMMA & Represents the comma symbol to separate parts \\
      PLUS & Represents the plus symbol to perform the addition\\
      MINUS & Represents the minus symbol to perform the subtraction\\
      TIMES & Represents the star symbol to perform multiplications\\
      DIV & Represents the division symbol to perform divisions\\
      POW & Represents the power symbol to perform the power\\
      DOT & Represents the dot symbol used for accessing object elements\\
      EQUAL & Represents the equal symbol used in conditional comparison\\
      LOWER & Represents the lower symbol used in conditional comparison\\
      LOWER\_EQ & Represents lower or equal symbol used in conditional comparisons\\
      ASSIGN & Represents the local variables assignments\\
      \hline
      \end{tabular}
      \caption{List of Tokens\label{tab:listToken}}
    \end{table}

\subsection{Syntax Analysis}
  \label{sec:syntax}
  We will build a parser for our language. To do that, we will use \href{https://www.gnu.org/software/bison/}{GNU Bison} 

  The parser is built from a grammar that is defined in the assignment.

  Each function is contained in the class. Each block contained in a function has as parent that particular function.


  For this part, we succeeded in all automated tests.

  \subsubsection{Building the tree}
    \label{sec:tree}
    The tree is composed of classes inheriting from AstNodes classes. The root node is a ProgramNode which contains a ClassNode for each class. 
    Each AstNode has a pointer pointing to their parent.

    This tree will be used differently in the following steps.
    %the printTree, then the second (only for types and checks good types)

    Each class is defined in a file given in the \texttt{nodes/} directory.


\subsection{Semantic checking}
  \label{sec:semantic}
  We implemented this step in three phases. 
  First we implemented the Classes checking, then the types checking while retrieving them. After that, we determined the non-redefinition of methods or variables.

  The implementation is done in the Semantic class which is in the \texttt{semantic/} folder.

  The errors are held in the structure SemErr and ExprType which carries the required information about semantic checking.

  We passed all the automated tests.

  \subsubsection{Classes checking}
    In this part, we will check: 
    \begin{itemize}
    \item Class redefinition
    \item Parent definition
    \item Inheritance cycles
    \item Main class and method prototype existence

    \end{itemize}
    To do so, we will use a 2D array that contains the Class ID (string), and the ClassNode associated to it. 
    This array is useful to know if this type has already been declared in the next class that will be analyzed. 
    know if all classes extending others are all declared.
      
  \subsubsection{Types retrieving and checking}
    The second time we go through the tree, we needed mechanisms to know the closest common parent that we use to determine the output 
    of a dynamic instantiation. 
    This step has the purpose of checking that all types are correctly implemented, and that all expressions have the correct type. This is done using the existing tree of nodes, instead of building a separate scope tree.

  \subsubsection{Redefinition}
    We have to determine if the  fields are not redefined. 
    The methods can be overwritten if they agree to certain conditions. We implemented this by going through the nodes
    and asking the parent classes (if needed) to find out if it is a redefinition.
    We also check the arguments of the methods during this step.



\subsection{Code Generation}
  \label{sec:codeGeneration}

  We learned that a LLVM library existed. However, it was quite difficult to handle so we decided to develop 
  our own version LlvmManager. At first, we were using a llvm tutorial which had the version 4 or 3.8. 
  However, we realized we needed to use version 3.5. This is why we used 

    \subsubsection{How it works}
    First, the LLVM will generate the header of the code. It writes down the methods types declarations. 

    Then, the LLVM will generate the class declaration.

    After that it declares the Main class and allocates the methods variables.

    It finishes by going through the tree to get the llvm translation of each node and it finishes by declaring the strings as global variables.
  \subsubsection{Choices and / or constraints}
    \begin{itemize}
    \item We saw that when we try to use \%0, llvm wants to use it for its own purpose (maybe in version 4). So we chose to begin
    giving names to our temporary variables beginning to \%1.
    \item We first implemented the instantiation of a new object as a block of operations allocating and initializing the objects. 
    At that time, we weren't really aware of the technique of using a knit method to do so. We then decided not to implement it
    since the main purpose of llvm code is not to have the greatest readibility. The only drawback of keeping or first idea is just making the llvm file a bit heavier
    \item We saw that getelemntptr in LLVM requires a load just afterwards. So we did it every time.
    \item The functions always takes as first argument the element on which they act. It allows us to use the keyword `self` accordingly.
    \end{itemize}

  \section{Code structure}
    This section is used to explain the structure of the code and its purpose.

\subsection{Doxygen}
  \label{Doxygen}
  We were able to generate Doxygen files. You are able to generate it using \texttt{make doc}. This command also displays the documentation directly in Firefox.
  Please have a look at \url{../doxygen/html/index.html} from the \texttt{report/} folder.
  \subsection{Folders}
    The table \ref{fig:folders} reflects the folders distribution.
    \begin{table}[H]
      \begin{tabular}{|c|p{10cm}|}
      \hline
      Folder & Description\\
      \hline\hline
      . &  Default folder containing the files for Flex, GNU Bison and LLvmManager class  \\\hline
      doxygen/ & The documentation (please see subsection \ref{Doxygen}) \\\hline
      nodes/ & The classes used to build the tree (cf. \ref{sec:tree})\\\hline
      report/ & Contains the report pdf\\\hline
      semantic/ & Contains the Semantic class and the different declarations required in step \ref{sec:semantic}\\\hline
      vsopl/ & Embeds the additional elements for code generation \\\hline
      \end{tabular}
      \caption{\label{fig:folders}Folders descriptions}
    \end{table}
    
\section{Improvements}
  \label{sec:improvements}
  %Le problème héritage d'un cycle
  \begin{itemize}
  \item Add an initial function : Init. It will make program heavier because it requires saving registries, ... 
  due to yet another function call. We had to know if the LLVM compiler would optimize it so that we do not 
  do something that will be removed just after.
  \item We have a problem with the inheritance of classes cycling. We tried to limit this problematic but we need to find
  a more elegant solution.
  \end{itemize}


\section{Conclusion}

  We know that our implementation is not perfect. In section \ref{sec:improvements}, you could find the improvements that we made.

\end{document}